\section{An integral CM-type criterion at coefficient two}
\label{sec:cm-type-unit-criterion}

Let \(p\) be an odd prime, let \(n=p-1=2h\), and choose a
primitive residue \(g\) modulo \(p\). Write
\[
 F=\mathbf Q(\zeta_p),\qquad
 G_p=\operatorname{Gal}(F/\mathbf Q)=\langle s\rangle,
 \qquad s(\zeta_p)=\zeta_p^g,\qquad c=s^h.
\]
For \(a\in\mathbf F_p^\times\), denote the corresponding automorphism
by \(\sigma_a\). The half-interval CM type and its inverse are
\begin{equation}
 \Phi=\sum_{a=1}^{h}\sigma_a,
 \qquad \Phi^*=\sum_{a=1}^{h}\sigma_a^{-1}.
                                                        \label{eq:cm-type-1}
\end{equation}
This section gives the integral criterion used in the simultaneous CM
normalization. The exact arithmetic image is proved in
Section~\ref{sec:prime-degree-local-cm}. Every local factor above two
is included in the coefficient algebra.

\subsection{The explicit group-ring factor}

Define a periodic sequence, with period \(n\), by
\[
 \varphi_k=
 \begin{cases}
 1,&\text{if the least positive residue of }g^{-k}\text{ is at most }h,\\
 0,&\text{otherwise}.
 \end{cases}
\]
Thus \(\varphi_{k+h}=1-\varphi_k\), and
\(\Phi^*=\sum_{k=0}^{n-1}\varphi_k s^k\). Put
\begin{equation}
 Q_p(s)=\sum_{k=0}^{h-1}(\varphi_k-\varphi_{k-1})s^k,
 \qquad
 R_p=\mathbf Z_2[G_p]/(1+c)
       =\mathbf Z_2[s]/(s^h+1).
                                                        \label{eq:cm-type-2}
\end{equation}
The index \(-1\) in this formula means \(n-1\).

\begin{proposition}\label{prop:explicit-cm-type-factor}
One has the exact identity
\begin{equation}
                 (s-1)\Phi^*=(c-1)Q_p
                  \quad\text{in }\mathbf Z[G_p].        \label{eq:cm-type-3}
\end{equation}
The class of \(Q_p\) modulo \(1+c\) is uniquely determined by \eqref{eq:cm-type-3}.
It is a unit of \(R_p\) if and only if
\begin{equation}
        \gcd\bigl(\overline Q_p(X),X^h+1\bigr)=1
                          \quad\text{in }\mathbf F_2[X].\label{eq:cm-type-4}
\end{equation}
If \(h=2^v m\) with \(m\) odd, one may replace \(X^h+1\)
in \eqref{eq:cm-type-4} by \(X^m-1\). Moreover
\begin{equation}
                         Q_p(1)=\pm1.                  \label{eq:cm-type-5}
\end{equation}
Consequently every prime with \(p-1\) a power of two satisfies
\eqref{eq:cm-type-4}. The unit condition is independent of the choice of primitive
generator of \(G_p\), of translation of the CM type, and of replacing
the type by its complement.
\end{proposition}

\begin{proof}
The coefficient of \(s^k\) in \((s-1)\Phi^*\) is
\(\varphi_{k-1}-\varphi_k\). The relation
\(\varphi_{k+h}=1-\varphi_k\) makes these coefficients
anti-invariant under \(c\). For \(0\leq k<h\), the coefficients
of \((c-1)Q_p\) at \(s^k,s^{k+h}\) are respectively
\(-\varphi_k+\varphi_{k-1}\) and
\(\varphi_k-\varphi_{k-1}\). This proves \eqref{eq:cm-type-3}.

For any integral cyclic group ring of even order, the image of
\(c-1\) is exactly its anti-invariant subgroup: an anti-invariant
element has coefficient pairs \((a_k,-a_k)\). The kernel of
\(c-1\) is exactly \((1+c)\mathbf Z[G_p]\), because its
coefficient pairs are equal. This proves uniqueness with the stated
modulus, without division by two.

The ring \(R_p\) is finite free over \(\mathbf Z_2\). An element
is invertible if and only if its reduction modulo two is invertible:
one direction is immediate, and in the other direction an inverse
modulo two is corrected by the convergent geometric series for an
element of \(1+2R_p\). This gives \eqref{eq:cm-type-4}. In characteristic two,
\(X^h+1=(X^m-1)^{2^v}\), proving the reduced test.
Telescoping and the half-period relation give
\[
 Q_p(1)=\varphi_{h-1}-\varphi_{n-1}
                    =2\varphi_{h-1}-1.
\]
This proves \eqref{eq:cm-type-5}, including the asserted case \(m=1\).

If a different generator is \(s^u\), take \(u>0\) coprime to
\(n\). Its factor in the same ring is
\[
       Q_p\,(1+s+\cdots+s^{u-1}).
\]
The second factor is already a unit in
\(\mathbf Z[s]/(s^h+1)\). Indeed its determinant as a
\(\mathbf Z\)-linear multiplication operator is
\[
 \prod_{\zeta^h=-1}\frac{\zeta^u-1}{\zeta-1}=1:
\]
the map \(\zeta\mapsto\zeta^u\) permutes these roots and none
of them is one. A determinant-one integral operator has an integral
inverse. Translation of the type multiplies \(Q_p\) by a group
element, and complementation changes its sign, because
\((s-1)\sum_{\gamma\in G_p}\gamma=0\). These operations preserve
the unit condition.
\end{proof}

For completeness, the half-interval type has trivial multiplicative
stabilizer. If \(a\{1,\ldots,h\}=\{1,\ldots,h\}\) modulo \(p\),
then summing gives
\[
 a\frac{h(h+1)}2=\frac{h(h+1)}2=-\frac18\neq0
                              \quad\text{in }\mathbf F_p.
\]
Hence \(a=1\). This is the group-theoretic assertion needed when
identifying the reflex field of this CM type with \(F\); it does
not by itself identify any arithmetic inertia image.

\subsection{The full local norm-one torus}

Set
\begin{equation}
 A=\mathcal O_F\otimes\mathbf Z_2,\qquad
 U=1+2A,\qquad
 T=\{z\in U:z\,c(z)=1\}.                                \label{eq:cm-type-6}
\end{equation}
The algebra \(A\) is a product of unramified discrete valuation
rings. Their residue multiplicative groups have odd order. Thus
\(U\) is the maximal pro-two subgroup of \(A^\times\), and
\(T\) is the maximal pro-two subgroup of its norm-one torus.
For a group-ring element, exponent notation denotes its natural
multiplicative action on these abelian groups; for example
\((s-1)z=s(z)/z\).

\begin{theorem}\label{thm:cm-type-torus-unit}
The action of \(\mathbf Z_2[G_p]\) on \(T\) factors through
\(R_p\). The following conditions are equivalent:
\begin{enumerate}
\item \(Q_p\) is a unit of \(R_p\);
\item \(Q_p:T\longrightarrow T\) is an automorphism;
\item \((s-1)\Phi^*(U)=T\).
\end{enumerate}
For every \(p\), without a unit hypothesis, the exact identity is
\begin{equation}
                     (s-1)\Phi^*(U)=Q_pT.               \label{eq:cm-type-7}
\end{equation}
If \(\mathcal D_p\) is the determinant of multiplication by
\(Q_p\) on the free rank-\(h\) module \(R_p\), then the
unit criterion is equivalently \(\mathcal D_p\) odd. In particular
it is an integral determinant criterion, with no class-number
assumption.
\end{theorem}

\begin{proof}
The relation \((1+c)T=0\) is the definition of \(T\). Continuous
\(\mathbf Z_2\)-exponents on a pro-two abelian group define the
claimed ring action. If \(Q_p\) is a unit, its inverse in \(R_p\)
acts as its continuous inverse on \(T\).

The relative unit version of Hilbert 90 gives
\begin{equation}
                             (c-1)U=T.                  \label{eq:cm-type-8}
\end{equation}
Here is the argument for the product algebra in \eqref{eq:cm-type-6}. If \(c\)
interchanges two field factors, a norm-one pair has the form
\((z,z^{-1})\), and is the quotient of a principal-unit pair
by its \(c\)-conjugate. If \(c\) fixes a factor, it acts as the
nontrivial automorphism of an unramified quadratic extension of
the fixed field. Ordinary Hilbert 90 writes a norm-one principal
unit as \(c(v)/v\). Removing a power of two makes \(v\) a unit.
Its residue is fixed by \(c\), since the quotient has residue one.
Dividing by the Teichmuller lift of that fixed residue makes
\(v\) a principal unit. This proves \eqref{eq:cm-type-8} in every orbit of
factors. For the ordinary Hilbert 90 and unramified unit norm
statements, see \cite[II, Proposition~1.22 and Corollary~1.23,
p.~67; III, Propositions~1.1--1.2, pp.~97--98]{MilneCFT}.
Combining \eqref{eq:cm-type-8} with \eqref{eq:cm-type-3} proves \eqref{eq:cm-type-7}.

It remains to justify the converse, including integral rather than
merely rational invertibility. The open subgroup
\(T(4)=T\cap(1+4A)\) is identified by the two-adic logarithm with
\begin{equation}
      4A^-,\qquad A^-=\{a\in A:c(a)=-a\}.               \label{eq:cm-type-9}
\end{equation}
The elements \(\zeta_p^a\), \(1\leq a\leq p-1\), form an
integral basis of \(\mathcal O_F\): they replace the element one
of the usual power basis using
\(1=-\sum_{a=1}^{p-1}\zeta_p^a\), with an integral inverse
change of basis. The group \(G_p\) permutes this basis regularly.
Consequently \(A\simeq\mathbf Z_2[G_p]\), and the elementary
anti-invariant calculation used in Proposition
\ref{prop:explicit-cm-type-factor} gives \(A^-\simeq R_p\).
Thus \eqref{eq:cm-type-9} is a free rank-one \(R_p\)-module.

The group \(T\) is a finitely generated \(\mathbf Z_2\)-module
with finite torsion, and its torsion-free quotient has the same
rational representation as \eqref{eq:cm-type-9}. If \(Q_p\) is onto \(T\),
it is onto that free quotient, so its determinant there is a
two-adic unit. That determinant is \(\mathcal D_p\), because
determinants are independent of the chosen integral lattice in
the rational representation. Multiplication by \(Q_p\) on
\(R_p\) is therefore invertible. This proves both the converse
and the equivalence with condition three.
\end{proof}

\subsection{The exact coinvariants and their quadratic characters}

For \(t\in\mathbf Z_2^\times\), put
\[
 a(t)=\frac{t^2-1}{8}\pmod2,
 \qquad b(t)=\frac{t-1}{2}\pmod2.
\]
These two characters identify
\(\mathbf Z_2^\times/(\mathbf Z_2^\times)^2\) with
\(\mathbf F_2^2\).

\begin{theorem}\label{thm:full-cyclotomic-torus-coinvariants}
There is a canonical isomorphism
\begin{equation}
 \Delta_s:\mathbf Z_2^\times/(\mathbf Z_2^\times)^2
       \xrightarrow{\ \sim\ }T/(s-1)T,
 \qquad
 t\longmapsto[s(v)/v],\quad v\in U,\quad v\,c(v)=t.
                                                        \label{eq:cm-type-10}
\end{equation}
In particular
\begin{equation}
 T^{G_p}=\{1,-1\},\qquad
 T/(s-1)T\simeq C_2^2,
 \qquad [T:(s-1)T]=4.                                  \label{eq:cm-type-11}
\end{equation}
Define
\(\chi_a=a\circ\Delta_s^{-1}\) and
\(\chi_b=b\circ\Delta_s^{-1}\) on \(T\). Then
\begin{equation}
             (s-1)T=\ker\chi_a\cap\ker\chi_b.           \label{eq:cm-type-12}
\end{equation}
They are exactly the two invariant torus characters of the
coefficient-two cyclic-cover construction: \(\chi_a\) defines
\(\xi_{2,p}\), and \(\chi_b\) defines its companion
\(\widetilde\psi_p\). Thus the companion equality
\(\widetilde\psi_p=\rho_p^I/2\bmod2\) identifies which known
Kummer coordinate is removed in \eqref{eq:cm-type-12}.
\end{theorem}

\begin{proof}
The relative norm \(U\to U^c\) is onto. For paired factors this
is immediate. For an internal unramified quadratic extension,
surjectivity of the unit norm follows first on residue units and
then on successive principal-unit quotients by the surjectivity
of the residue trace. If the target is a principal unit, the
residue of an initial lift has norm one, and dividing by its
Teichmuller lift gives a principal-unit lift. This proves the
existence of \(v\) in \eqref{eq:cm-type-10}.

The quotient \(s(v)/v\) has relative norm one, because \(t\)
is rational. Replacing \(v\) by another norm lift changes this
quotient by an element of \((s-1)T\). A rational square \(t=w^2\)
has the lift \(v=w\), so the indicated map is well defined on
square classes and is a homomorphism.

We spell out the Hilbert 90 step for the full product algebra.
The algebra \(E=F\otimes\mathbf Q_2\), with its \(G_p\)-action,
is a cyclic Galois algebra over \(\mathbf Q_2\). Its factors are
permuted transitively by \(s\); the stabilizer of a factor is the
decomposition group \(\langle\sigma_2\rangle\). Cyclic Hilbert
90 for this algebra follows directly by solving the equations
successively around the cycle of factors and then applying
ordinary Hilbert 90 to the resulting equation in one field
factor. The consistency condition is precisely the full norm
equal to one. This argument does not require \(E\) to be a field.

For \(z\in T\), the full norm to \(\mathbf Q_2\) is one,
since \(N_{G_p}=(1+s+\cdots+s^{h-1})(1+c)\). Thus
\(z=s(v)/v\) for \(v\in E^\times\). Since \(z\) is a unit,
the valuation vector of \(v\) is \(G_p\)-fixed. Transitivity
and unramifiedness make all valuations equal integers; dividing
by one rational power of two makes \(v\) a unit in every factor.
Since \(z\equiv1\pmod{2A}\), its residue satisfies \(s(v)=v\).
The invariant subalgebra of \(A/2A\) is \(\mathbf F_2\):
invariants in one factor are the Frobenius-fixed elements, and
the other factors are determined by transitivity. Its only
nonzero unit is one. Consequently \(v\in U\).

Now \(s(v\,c(v))=v\,c(v)\), so its relative norm is in
\(U^{G_p}=\mathbf Z_2^\times\). This proves surjectivity of
\eqref{eq:cm-type-10}. If \(s(v)/v=s(z)/z\) for some \(z\in T\), then
\(v/z\) is a rational unit, and \(v\,c(v)\) is its square.
This proves injectivity. All the groups are compact and the
images of the continuous maps are closed, so the isomorphism
is continuous. The fixed subgroup assertion follows from
\(U^{G_p}=\mathbf Z_2^\times\): relative norm one means \(t^2=1\).
This proves \eqref{eq:cm-type-11} and \eqref{eq:cm-type-12}.

We compare the characters with their local-factor definitions.
Let \(d=\operatorname{ord}_p(2)\). If \(d=2e\) is even,
there are \(r=n/d\) self-dual factors, each an unramified
extension \(L/\mathbf Q_2\) of degree \(2e\), and
\(c\) acts there as \(\sigma^e\), with \(\sigma\) Frobenius.
The local norm-one coinvariant map is
\[
 \Delta_{\sigma}(t)=[\sigma(v)/v],\qquad v\sigma^e(v)=t.
\]
The established local characters are the sums, over the \(r\)
factors, of \(a\circ\Delta_{\sigma}^{-1}\) and
\(b\circ\Delta_{\sigma}^{-1}\). Choose factor representatives
\(1,s,\ldots,s^{r-1}\). The last passage around them is
\(\sigma^k\), where \(s^r=\sigma_2^k\) and
\(\gcd(k,d)=1\); in particular \(k\) is odd. Choose the same
norm lift of \(t\) in these coordinates. All quotients cancel
except the last, \(\sigma^k(v)/v\). Its local coinvariant is
\(\Delta_{\sigma}(t)^k=\Delta_{\sigma}(t)\). Thus the two
sums take the values \(a(t),b(t)\) on \eqref{eq:cm-type-10}.

If \(d\) is odd, the factors occur in \(r=n/(2d)\) reciprocal
pairs. Write their norm-one coordinates as \((z_i,z_i^{-1})\)
with \(z_i\) in an unramified extension of degree \(d\).
The established characters are
\[
        \sum_i a(N_{L/\mathbf Q_2}z_i),\qquad
        \sum_i b(N_{L/\mathbf Q_2}z_i).
\]
Choose the norm lift \((1,t)\) on every pair. A primitive
generator cycles the pairs, and the last passage exchanges
the reciprocal factors, possibly with a Frobenius twist.
Its quotient has a single torus coordinate \(t\) or \(t^{-1}\).
The norm is \(t^{\pm d}\); since \(d\) is odd, its two
characters are again \(a(t),b(t)\). This proves the claimed
identification in both cases.
\end{proof}

The group \((s-1)T\) has \(\mathbf Z_2\)-rank \(h\), because
it is open in \(T\). It need not be torsion-free: when the local
algebra has several orbits under complex conjugation, differences
of their signs can give nontrivial torsion in this image. No
identification of this entire compact group with a free
\(\mathbf Z_2\)-module is intended.

In particular, if an arithmetic construction establishes a first
commutator image equal to \(T\), its further commutator by an
element acting as \(s\) has exact image \((s-1)T\). Formula
\eqref{eq:cm-type-12} describes the two quadratic torus conditions on that
second image. This implication is purely algebraic; it does not
assert that a single inertia group supplies all the local factors.

\subsection{Exact indices and illustrative primes}

There is also a finite formula for the possible failure of the
unit condition. Put
\[
 d=\operatorname{ord}_p(2),\qquad
 r=\begin{cases}
       n/d,& d\text{ even},\\
       n/(2d),&d\text{ odd}.
    \end{cases}
\]
Each unramified factor has only \(\{1,-1\}\) as pro-two
torsion. The relative norm identifies the two signs in a paired
orbit. Thus
\begin{equation}
 T_{\mathrm{tors}}\simeq(\mathbf F_2)^r
              \simeq\mathbf F_2[s]/(s^r-1),              \label{eq:cm-type-13}
\end{equation}
where \(s\) permutes the \(r\) orbits of complex conjugation
transitively. Define
\[
 e_p=\deg\gcd(\overline Q_p(X),X^r-1).
\]
Then \(\ker(Q_p|T_{\mathrm{tors}})\) has order \(2^{e_p}\).
If \(\mathcal D_p\neq0\), the exact image index is
\begin{equation}
                  [T:Q_pT]=2^{v_2(\mathcal D_p)+e_p}.    \label{eq:cm-type-14}
\end{equation}
Indeed, on the free quotient of \(T\) the kernel is zero and
the cokernel has order \(2^{v_2(\mathcal D_p)}\). The exact
sequence for the torsion subgroup and its free quotient then
multiplies this index by the order of the torsion cokernel,
which equals the order of its kernel. If \(\mathcal D_p=0\),
the image has infinite index. No universal nonvanishing of this
determinant is needed for the criterion or asserted here.

Using the least primitive positive residue as \(g\), the first
examples are as follows. Each displayed determinant is the
integer resultant \(\operatorname{Res}_X(X^h+1,Q_p(X))\).
\[
\begin{array}{c|c|l|r|c}
 p&g&Q_p(X)&\mathcal D_p&[T:Q_pT]\\\hline
 5&2&-X&1&1\\
 7&3&-X&1&1\\
 11&2&-X+X^2-X^3&3&1\\
 13&2&-X+X^3-X^4&5&1\\
 17&3&-X^3+X^4-X^7&17&1\\
 19&2&-X+X^2-X^3+X^4-X^6&27&1\\
 23&5&-X+X^3-X^5+X^6-X^7&267&1\\
 29&2&-X+X^3-X^4+X^5-X^7+X^{10}-X^{11}&4520&8
\end{array}
\]
For \(p=29\), the obstruction is visible without a large
determinant:
\[
 \gcd(\overline Q_{29}(X),X^{14}+1)=X^3+X^2+1.
\]
Here \(d=28\), so \(r=1\), and \(Q_{29}(1)=-1\) gives
\(e_{29}=0\). Since \(4520=8\cdot565\), \eqref{eq:cm-type-14} gives the
stated exact index. Thus the unit criterion excludes
some primes. Further exact failures occur at \(113\), with
the same cubic greatest common divisor, and at \(163\), with
greatest common divisor \(X^2+X+1\); their indices are respectively
eight and four. These are finite arithmetic computations of
\eqref{eq:cm-type-2}--\eqref{eq:cm-type-4}, not an extrapolation or a class-number argument.

The case \(p=7\) illustrates why the product algebra matters.
Here \(d=3\), so \(A\) has two unramified cubic factors and
complex conjugation interchanges them. Nevertheless \(Q_7=-s\)
is an integral unit, \(T\) is the principal-unit group of one
cubic factor, and \eqref{eq:cm-type-10}--\eqref{eq:cm-type-12} still give the full two-dimensional
quadratic quotient. Establishing that the relevant arithmetic
inertia groups realize this full product-unit calculation is a
separate CM reciprocity step.

\begin{proposition}\label{prop:cm-commutator-torus-torsion}
With the notation of Theorem
\ref{thm:full-cyclotomic-torus-coinvariants}, let \(P=(s-1)T\)
and let \(r\) be the number of orbits of complex conjugation on
the primes above two. Then
\[
 P_{\mathrm{tors}}=(s-1)T_{\mathrm{tors}}
       \simeq C_2^{r-1}.
\]
In particular, as compact abelian groups, noncanonically,
\[
       T\simeq\mathbf Z_2^h\times C_2^r,
       \qquad P\simeq\mathbf Z_2^h\times C_2^{r-1}.
\]
For example, \(p=17\) has \(r=2\), so \(P\) has one
nontrivial element of order two, even though the unit criterion
holds.
\end{proposition}

\begin{proof}
Under \eqref{eq:cm-type-13}, the image \((s-1)T_{\mathrm{tors}}\) is the
even-parity sign subgroup and has dimension \(r-1\) over
\(\mathbf F_2\). We show that the restriction of
\((\chi_a,\chi_b)\) to \(T_{\mathrm{tors}}\) has kernel
exactly this subgroup. This will prove equality with
\(P\cap T_{\mathrm{tors}}\), which is the torsion subgroup
of \(P\).

If \(d=\operatorname{ord}_p(2)\) is odd, a sign in one paired
factor has odd-degree norm \(-1\). Its character values are
\(a(-1)=0\) and \(b(-1)=1\). Thus the restriction is
\((0,\sum_i\epsilon_i)\) on the sign vector
\((\epsilon_1,\ldots,\epsilon_r)\).

If \(d=2e\) is even, let \(\sigma\) be Frobenius in one
unramified factor. Its unramified quadratic subfield contains
\(\delta=1+2\omega\), with \(\omega^2+\omega+1=0\).
Then \(\delta^2=-3\), \(\sigma(\delta)=-\delta\), and
\[
 \delta\,\sigma^e(\delta)=
 \begin{cases}3,&e\text{ odd},\\-3,&e\text{ even}.
 \end{cases}
\]
Hence the local coinvariant of \(-1\) corresponds to the
square class of \(3\) or \(-3\), respectively. Its two
character values are \((1,e\bmod2)\). On all factors the
restriction is therefore
\((\sum_i\epsilon_i,e\sum_i\epsilon_i)\). In both cases
its kernel is the even-parity sign subgroup, as required.

Finally, a finitely generated compact \(\mathbf Z_2\)-module
splits as its finite torsion subgroup and a free module.
Both \(T\) and its open subgroup \(P\) have rank \(h\),
which gives the two displayed decompositions.
\end{proof}
